\section{Architettura della Pipeline}
\label{sec:architettura}

L'architettura sviluppata per questo elaborato orchestra il processo di inferenza, estrazione vettoriale, riduzione dimensionale e salvataggio dei dati. 

L'intera logica di esecuzione è stata implementata all'interno del notebook Python \texttt{llm-semantic-explorer-pipeline.ipynb} \cite{notebook_pipeline}. 

Di seguito vengono analizzate le singole fasi del flusso di elaborazione.

\subsection{Generazione e Filtraggio dei Dati}
\label{sec:gestione_dati}

La prima fase della pipeline riguarda l'interazione con il modello linguistico e la gestione dei prompt. Per poter misurare l'effetto dello \textit{steering} semantico, il sistema è configurato per processare iterativamente coppie o catene di prompt. Ad ogni tecnica avanzata (come \textit{Chain of Thought}, \textit{Few-Shot} o \textit{Persona Prompting}) viene affiancato un prompt ``di base'' (\textit{zero-shot}) che funge da traiettoria di controllo.

Per isolare l'impatto della struttura del prompt sul ragionamento ed evitare rumore generato da convenzioni colloquiali, al modello viene fornito un \textit{system prompt} restrittivo. Tale istruzione di sistema impone la generazione di risposte dirette e prive di preamboli (es. formule di saluto o conclusioni discorsive), forzando l'LLM ad emettere esclusivamente i token strettamente necessari all'esecuzione del task.

Una volta avviato il processo di inferenza, il modello genera sequenzialmente i token che compongono la risposta. Tuttavia, la proiezione dell'intero output produrrebbe una mappa tridimensionale eccessivamente densa di elementi puramente grammaticali o sintattici, privi di reale rilevanza topologica. Per ovviare a questo problema, la pipeline implementa una routine di filtraggio post-generazione:
\begin{itemize}
    \item \textbf{Rimozione della punteggiatura e dei caratteri speciali}: vengono scartati i token contenenti esclusivamente simboli, spaziature isolate o tag di sistema (come i marcatori di fine sequenza).
    \item \textbf{Rimozione delle \textit{stop-words}}: vengono eliminati articoli, preposizioni e congiunzioni avvalendosi del vocabolario fornito dalla libreria NLTK.
\end{itemize}

L'applicazione di questa funzione di filtraggio fa sì che all'interno della struttura dati vengano conservati unicamente i token portatori di un significato semantico definito. Solo per questo sottoinsieme filtrato verranno mantenuti in memoria i corrispondenti vettori ad alta dimensionalità, ottimizzando così il carico computazionale per le successive fasi di clustering e riduzione.

\subsection{Estrazione degli Hidden States dai Layer del Modello}
\label{sec:estrazione_hidden_states}

La fase centrale dell'architettura sperimentale consiste nell'estrazione delle rappresentazioni vettoriali durante il processo generativo del modello. Poiché l'obiettivo dell'elaborato è misurare lo spazio latente in modo dinamico, la semplice funzione \texttt{generate} non è sufficiente, in quanto di default restituisce esclusivamente la sequenza di identificativi (\textit{token IDs}) che costituiscono il testo di output.

Per catturare le attivazioni interne della rete neurale (\textit{hidden states}), la pipeline implementa un passaggio aggiuntivo (\textit{forward pass}) eseguito parallelamente alla generazione standard. Sfruttando le specifiche fornite dalla libreria \textit{Transformers}, è possibile inviare al modello una richiesta di processamento abilitando il parametro \texttt{output\_hidden\_states} \cite{huggingface_transformers_docs}. Questa impostazione istruisce la rete a restituire non solo la predizione del vocabolario, ma anche la matrice multidimensionale prodotta in uscita da ogni singolo layer.

Al fine di ottenere una rappresentazione sintetica ma semanticamente completa, il sistema intercetta unicamente i tensori in uscita dall'ultimo layer della rete (accessibili all'indice -1 dell'array restituito dal modello). Come consolidato in letteratura \cite{belrose2023eliciting}, l'ultimo strato di una rete auto-regressiva contiene infatti i vettori contestuali definitivi, i quali aggregano tutte le informazioni sintattiche e semantiche elaborate dai layer precedenti e vengono impiegati direttamente per la proiezione finale.

Dal punto di vista dell'implementazione tecnica, l'estrazione avviene all'interno dell'ambiente \textit{PyTorch} associato all'acceleratore hardware. I vettori estratti assumono la forma di tensori ad alta dimensionalità allocati nella memoria video (VRAM). Per rendere questi dati compatibili con i successivi algoritmi di analisi geometrica, i tensori vengono trasferiti dalla VRAM alla memoria di sistema della CPU e successivamente convertiti dal formato tensoriale \textit{PyTorch} al formato matriciale \textit{NumPy}. Questa matrice costituisce il dataset matematico su cui verranno eseguiti la riduzione dimensionale e il calcolo delle metriche.

\subsection{Integrazione dei Processi di Riduzione Dimensionale e Clustering}
\label{sec:integrazione_riduzione_clustering}

Ottenuta la matrice filtrata degli ultimi \textit{hidden states}, il sistema procede all'estrazione della topologia spaziale. La gestione di questo processo si articola in tre stadi computazionali, progettati per aggirare il fenomeno della \textit{Curse of Dimensionality} \cite{bellman1961adaptive}, critico per gli algoritmi basati sulla distanza in spazi ad alta dimensionalità \cite{altman2018curse}.

Il primo stadio consiste in una riduzione dimensionale propedeutica, applicando l'algoritmo UMAP per comprimere i vettori originali da 3072 a 30 dimensioni. Questo valore rappresenta un compromesso euristico ottimale: è sufficientemente basso da ripristinare il potere discriminante della metrica di distanza spaziale, ma abbastanza elevato da trattenere l'intricata informazione semantica necessaria al partizionamento. 

Il secondo stadio prevede l'applicazione dell'algoritmo HDBSCAN sui vettori a 30 dimensioni. L'algoritmo è stato configurato per adattarsi alla specifica natura linguistica dei dati:
\begin{itemize}
    \item \texttt{min\_samples = 2}: regola la fase di esplorazione locale definendo il numero minimo di vicini affinché un punto formi un nucleo (\textit{core point}). Un valore così basso abbassa la soglia di tolleranza, permettendo di preservare e raggruppare anche i \textit{token} che formano le frasi di transizione logica, essenziali per tracciare la continuità del ragionamento del modello.
    \item \texttt{min\_cluster\_size = 5}: regola la soglia di sopravvivenza finale dei gruppi. Essendo ogni vettore un singolo \textit{token}, un limite di 5 unità garantisce l'identificazione di micro-concetti validi (es. proposizioni chiave di poche parole), impedendo al contempo la formazione di aggregazioni puramente accidentali.
\end{itemize}

A valle del clustering vettoriale, la pipeline esegue un processo di \textit{Cluster Labeling} automatizzato. Per ogni gruppo identificato da HDBSCAN, l'insieme dei token testuali corrispondenti viene raggruppato e fornito in input allo stesso LLM. Tramite un prompt di sistema restrittivo (che vincola l'output a un massimo di 1-5 parole), il modello genera un'etichetta semantica che descrive sinteticamente l'argomento trattato in quella specifica regione dello spazio latente. Per garantire che l'LLM operi in maniera più  rigorosa, il parametro \textit{temperature} in questa fase è stato abbassato a 0.1, forzando un comportamento strettamente deterministico che massimizza la riproducibilità e l'oggettività dell'etichetta generata.

Il terzo e ultimo stadio riguarda la proiezione spaziale definitiva per l'esportazione. Una seconda istanza indipendente dell'algoritmo UMAP viene applicata alla medesima matrice di \textit{hidden states} originaria per comprimerla in uno spazio visibile (\texttt{n\_components = 3}). Per preservare la continuità visiva della traiettoria di generazione, i parametri vengono configurati specificamente \cite{mcinnes2018umap}: \texttt{n\_neighbors} è fissato a 15, bilanciando la struttura locale (token adiacenti nella frase) con quella globale (paragrafi diversi), mentre \texttt{min\_dist} è fissato a 0.1, garantendo che i token semanticamente identici mantengano una minima separazione fisica, evitando l'eccessivo collasso visivo dei punti e consentendone l'ispezione interattiva.

L'output finale, composto da coordinate tridimensionali, etichette matematiche HDBSCAN e tag semantici, viene aggregato e serializzato in formato JSON, strutturando la base di dati per l'analisi quantitativa e la navigazione 3D.